\section{Results}
\label{sec:results}

This section reports speculative-decoding outcomes from
the consolidated 12-configuration summary CSV and baseline context from the
deterministic and stochastic baseline CSV artifacts. We emphasise paired comparative
metrics ($S$, $\alpha$, $B_{\text{eff}}$) and use task deltas to evaluate the
speed--quality tradeoff.

\begin{table*}[t]
\centering
\caption{Speculative grid summary (Qwen2.5-3B target). $S$ is paired speedup
reported by the summary artifact; $\alpha$ is token-level acceptance;
$B_{\text{eff}}$ is accepted tokens per verification cycle.}
\label{tab:headline}
\small
\begin{tabular}{l c c c c c c c}
\hline
\textbf{Config} & $S$ & $\alpha$ & $B_{\text{eff}}$ & $\Delta$GSM8K & $\Delta$MMLU & $\Delta$CNN/DM & $R_{\text{tok}}$ \\
\hline
0.5B, $k{=}4$, det   & 1.6565 & 0.4168 & 1.65 & -1.67 & 0.60 & -0.09 & 12.37 \\
0.5B, $k{=}4$, stoch & 1.5706 & 0.4058 & 1.60 & 5.34  & 0.60 & -0.11 & 12.00 \\
0.5B, $k{=}8$, det   & 1.7130 & 0.3412 & 2.64 & -0.67 & 0.20 & -0.43 & 12.94 \\
0.5B, $k{=}8$, stoch & 1.6128 & 0.3380 & 2.62 & 2.00  & 0.20 & -0.73 & 12.25 \\
0.5B, $k{=}16$, det  & 1.4629 & 0.2410 & 3.57 & -2.00 & 1.40 & -0.19 & 11.74 \\
0.5B, $k{=}16$, stoch& 1.4059 & 0.2355 & 3.50 & 3.00  & 0.00 & -0.50 & 11.51 \\
1.5B, $k{=}4$, det   & 1.6504 & 0.4555 & 1.80 & -0.33 & -0.20 & -0.17 & 12.19 \\
1.5B, $k{=}4$, stoch & 1.5812 & 0.4401 & 1.74 & 2.67  & 0.20 & -0.05 & 12.10 \\
1.5B, $k{=}8$, det   & \textbf{1.7289} & 0.3829 & 2.96 & -1.00 & 0.60 & -0.37 & 12.80 \\
1.5B, $k{=}8$, stoch & 1.6215 & 0.3709 & 2.87 & 6.34  & 1.00 & -0.10 & 12.27 \\
1.5B, $k{=}16$, det  & 1.4968 & 0.2821 & 4.19 & 2.67  & -0.20 & -0.11 & 11.60 \\
1.5B, $k{=}16$, stoch& 1.4432 & 0.2781 & 4.13 & 4.00  & -1.00 & -0.50 & 11.55 \\
\hline
\end{tabular}
\end{table*}

\subsection{RQ1: Does speculative decoding accelerate inference?}

All 12 speculative configurations exceed $S{=}1$, with observed speedups
from $1.4059$ to $1.7289$. The best configuration is 1.5B with $k{=}8$ in
deterministic mode. Baseline context from separate regime runs shows mean
latencies of 11.41s (deterministic) and 11.65s (stochastic) per sample, so the
measured speculative gains are substantial in wall-clock terms for fixed
workload and prompts.

\subsection{RQ2: How do draft length and draft size affect acceptance and speed?}

Draft length exhibits a non-monotonic effect on speedup. Averaging over drafts
and regimes, mean speedup is highest at $k{=}8$ ($1.6690$), compared with
$k{=}4$ ($1.6147$) and $k{=}16$ ($1.4522$). This pattern is consistent with the
acceptance dynamics: as $k$ increases, $B_{\text{eff}}$ increases but
$\alpha$ decreases, and beyond $k{=}8$ the additional drafting/verification cost
dominates net benefit.

Draft size has a smaller global effect than $k$ in this grid
(mean $S$: 1.5703 for 0.5B versus 1.5870 for 1.5B), but the larger draft is
advantageous at the best operating point ($k{=}8$, deterministic).

\subsection{RQ3: Deterministic versus stochastic regime}

Deterministic decoding is consistently faster for every matched draft--$k$
pair. The deterministic advantage ranges from $0.0536$ to $0.1074$ in absolute
$S$ and averages $0.0789$ across the six matched pairs. This supports a
deployment preference for deterministic regime when latency is primary and
output diversity is not required.

Quality deltas show different stability by task. CNN/DailyMail remains close to
baseline across all settings (from $-0.73$ to $-0.05$). MMLU variation is also
limited (from $-1.0$ to $1.4$). GSM8K is the most sensitive axis, especially
under stochastic settings where positive and negative swings are larger. Under
a conservative quality filter ($\max |\Delta| \le 1$ across tasks), the best
configuration remains 1.5B, $k{=}8$, deterministic.

\subsection{Summary of empirical operating point}

The current evidence supports a practical recommendation: use deterministic
speculative decoding with $k{=}8$ and prefer the 1.5B draft when memory budget
permits; otherwise 0.5B at $k{=}8$ is a strong alternative. Increasing to
$k{=}16$ raises accepted block size but lowers net speedup in this setup,
indicating that acceptance gains no longer offset added cycle cost.
